\documentclass[12pt]{article}
\usepackage[utf8]{inputenc}
\usepackage[T1]{fontenc}
\usepackage[margin=1in]{geometry}
\usepackage{setspace}
\usepackage{xcolor}
\usepackage[colorlinks=true, urlcolor=blue!60!black, linkcolor=blue!60!black]{hyperref}
\setlength{\parindent}{0em}
\setlength{\parskip}{0.6em}
\renewcommand{\baselinestretch}{1.0}
\begin{document}
\begin{center}
{\LARGE \textit{Gundy v. United States}, 588 U.S. 128 (2019)\par}
\vspace{0.5em}
{\large Case Brief --- Two-Page Edition\par}
\end{center}

\section{Memory Jogger}
SORNA's delegation to the Attorney General to apply the Act to pre-Act sex offenders survived: read to require registration ``as soon as feasible,'' it had an intelligible principle (\href{https://supreme.justia.com/cases/federal/us/588/17-6086/}{Justia}).

\section{Detailed Case Facts}
SORNA (2006) created a comprehensive national sex-offender registry and let the Attorney General ``specify the applicability'' of its requirements to some 500,000 offenders convicted before the Act. The Attorney General applied it to all of them. Herman Gundy, convicted in Maryland in 2005, moved to New York after his release without registering.

\section{Procedural History}
Convicted of failing to register, Gundy argued that \S~20913(d) was an unconstitutional delegation; the district court and the Second Circuit rejected the claim.

\section{Judicial Votes}
Affirmed, 5--3. Kagan wrote for Ginsburg, Breyer, and Sotomayor; Alito concurred in the judgment; Gorsuch, with Roberts and Thomas, dissented; Kavanaugh did not participate.

\section{Holding}
Section 20913(d) is constitutional. Construed to require registration of all pre-Act offenders as soon as feasible, it is a narrow transitional delegation that satisfies the intelligible-principle test.

\section{Analysis of Opinions}
\textbf{Plurality:} Statutory meaning comes first. The Act's comprehensive purpose, its past-tense definition of ``sex offender,'' its history, and its reference to offenders unable to comply confine the Attorney General to feasibility. So read, the delegation is ``distinctly small-bore'' compared with ``public interest'' standards long upheld. \textbf{Alito:} he would reconsider the doctrine with a willing majority, but not single out this law. \textbf{Gorsuch:} Congress may let the Executive fill up details, find facts, or exercise its own powers---not write a criminal code; the plurality rewrote the text. In 2025, \textit{FCC v. Consumers' Research} reaffirmed the intelligible-principle test.

\section{Examples: Future Applications}
\textbf{Same-Side:} a transitional FAA power to bring pre-statute drone operators into a registry is valid. \textbf{Opposite-Side:} a statute letting the Attorney General decide in her sole discretion whether a new felony applies to anyone likely falls. \textbf{Fence-Sitter:} a Commerce Secretary empowered to create felonies for frontier AI in the public interest and for national security.

\section{Critique}
Gorsuch charges that the Court saves broad statutes by rewriting them, and \textit{Gundy} applied the same lenient test to a felony as to civil ratemaking. Unable to draw a nondelegation line, the Court turned instead to removal power and the major questions doctrine.

\section{Key Quotations}
\begin{quote}
``Given that standard, a nondelegation inquiry always begins (and often almost ends) with statutory interpretation. Only after a court has determined a challenged statute's meaning can it decide whether the law sufficiently guides executive discretion to accord with Article I.''
\end{quote}
\noindent --- Justice Kagan, plurality opinion, 588 U.S. at 135--36.

\end{document}
